% ==================== 附录 ====================
\appendix
\renewcommand{\thesection}{附录 \Alph{section}}
\section{支撑材料文件列表}

本次提交的支撑材料为文件夹 \texttt{final\_version/}（压缩为单个文件提交），
其内容与下表的对应关系如下。
所有源程序均为完整、可直接运行版本，运行环境见 \S\ref{sec:appendix-env}。

\begin{table}[H]
\centering
\footnotesize
\setlength{\tabcolsep}{4pt}
\caption{支撑材料文件列表（\texttt{final\_version/} 目录结构）}
\label{tab:appendix-files}
\begin{tabularx}{\textwidth}{@{}l X c@{}}
\toprule
路径 & 说明 & 对应问题 \\
\midrule
\multicolumn{3}{@{}l}{\textbf{一、源程序}} \\
\midrule
\texttt{Q1/q1\_code/}      & \texttt{q1\_model.py} 数据读取与 LP/MILP 建模、约束核验、对偶量；\texttt{run\_q1.py} 主流程与导出；\texttt{plot\_q1.py} 绘图 & 问题一 \\
\midrule
\texttt{Q2/code/}          & \texttt{rolling\_window.py} 有效残差与候选窗口；\texttt{run\_experiment.py} 滚动执行；\texttt{verify\_window.py}、\texttt{check\_information.py} 核验；\texttt{report\_experiment.py} 汇总发布；\texttt{make\_target\_tables.py} 题目要求表 & 问题二 \\
\midrule
\texttt{Q3/code/}          & \texttt{q3\_data.py} 输入校验；\texttt{q3\_forecast.py} 预报融合；\texttt{q3\_model.py} 场景树 LP；\texttt{run\_q3.py} 滚动执行；\texttt{verify\_q3.py} 核验；\texttt{plot\_q3.py} 绘图 & 问题三 \\
\midrule
\texttt{Q4/part2/code/}    & \texttt{q42\_model.py} 电价预测与联合场景；\texttt{run\_q42.py} 主流程；\texttt{report\_q42.py} 汇总；\texttt{verify\_q42.py}、\texttt{check\_information.py} 核验 & 问题四-2 \\
\midrule
\texttt{Q4/part3/code/}    & 在问题三代码基础上增加实时电价场景与结算；\texttt{requirements.txt} 依赖清单 & 问题四-3 \\
\midrule
\texttt{common/q2\_base/}  & 第二问与 4-2 共用的基础 LP 与数据层（\texttt{q2\_data.py}、\texttt{run\_1439.py}） & 二 / 四 \\
\texttt{common/efficiency/} & \texttt{template\_layout.py} 附件5 结果模板表头（主程序读取模板所必需） & 一 / 二 / 三 / 四 \\
\texttt{common/plotting/}  & \texttt{setup\_style.py}、\texttt{export\_figure.py} 统一绘图样式与导出 & 全部 \\
\midrule
\multicolumn{3}{@{}l}{\textbf{二、结果与图表}} \\
\midrule
\texttt{Q1..Q4/*/results/} & \texttt{result*.xlsx} 正式结果（题目要求格式）、\texttt{summary.json} 汇总指标、\texttt{schedule\_detail.csv} 逐时段明细、\texttt{verification.json} 核验记录、\texttt{figures/} 插图源文件 & 全部 \\
\texttt{Q2/results/variants/}、\texttt{Q2/results/extensions/}、\texttt{Q2/results/sensitivity/} & 各候选窗口方案、短窗口扩展与敏感性分析的完整对照结果 & 问题二 \\
\texttt{Q3/results/comparisons/} & 四种发布时点方案（\texttt{B\_aligned}、\texttt{only\_0}、\texttt{at\_0\_6}、\texttt{at\_0\_6\_12}）的完整结果 & 问题三 \\
\texttt{Q4/part2/results/}、\texttt{Q4/part3/results/} & 各电价场景方案与指定日的完整结果 & 问题四 \\
\midrule
\multicolumn{3}{@{}l}{\textbf{三、数据}} \\
\midrule
\texttt{data/附件/}        & \texttt{附件1--4.xlsx} 原始数据；\texttt{附件5/} 结果文件模板（\texttt{result1.xlsx} 等 5 个） & 全部 \\
\bottomrule
\end{tabularx}
\end{table}

\textbf{说明}：\texttt{final\_version/} 为自包含目录，解压后无需任何外部文件
即可按 \S\ref{sec:appendix-env} 的命令复现全部结果；
附录 \S\ref{sec:appendix-source} 的源程序\textbf{直接取自该目录}，
与所提交文件逐字一致。

\section{运行环境与复现方法}
\label{sec:appendix-env}

\begin{table}[H]
\centering
\small
\caption{软件环境}
\begin{tabularx}{\textwidth}{@{}l X@{}}
\toprule
项目 & 版本 / 说明 \\
\midrule
操作系统 & Windows \\
Python   & 3.12 / 3.13 \\
数值计算 & \texttt{numpy}、\texttt{scipy}（HiGHS 求解器）\\
数据处理 & \texttt{pandas}、\texttt{openpyxl}\\
绘图     & \texttt{matplotlib} \\
论文排版 & XeLaTeX $+$ \texttt{ctex}（中文支持）\\
\bottomrule
\end{tabularx}
\end{table}

\textbf{复现步骤}（在 \texttt{final\_version/} 目录下执行）：

\begin{verbatim}
python -m pip install -r Q4/part3/requirements.txt

# 问题一
python -X utf8 Q1/q1_code/run_q1.py

# 问题二（推荐方案：有效残差 + 56 天等权）
python -X utf8 Q2/code/run_experiment.py \
       --first-residual 7 --only uniform56 --out recommended_only
python -X utf8 Q2/code/make_target_tables.py     # 题目要求的表1/表2/表3

# 问题三（四次更新主方案）
python -X utf8 Q3/code/run_q3.py --mode all --out reproduced
python -X utf8 Q3/code/verify_q3.py              # 独立核验

# 问题四
python -X utf8 Q4/part2/code/run_q42.py --primary joint_price
python -X utf8 Q4/part3/code/run_q3.py --mode all --out out
python -X utf8 Q4/part3/code/verify_q3.py        # 独立核验
\end{verbatim}

\textbf{说明}：
\begin{itemize}[itemsep=2pt,topsep=3pt]
  \item 各脚本内部路径均相对于文件位置，从任意工作目录用脚本绝对路径启动均可。
  \item 问题二\textbf{无参数默认运行的是原始残差口径候选}，
        不是最终推荐方案，须按上述命令显式指定 \texttt{--first-residual 7}。
  \item 问题四-2 会读取 \texttt{Q2/results/} 下的 Q2 结果作为同口径基准，
        因此需先运行问题二。
  \item 运行主程序会覆盖同名结果文件；论文使用的现有结果应先备份。
\end{itemize}

\section{核验记录摘要}

\begin{table}[H]
\centering
\caption{逐时段核验项目（四问合计 $48\,096$ 段/年）}
\small
\begin{tabular}{@{}l l c@{}}
\toprule
核验项 & 判据 & 结果 \\
\midrule
能量平衡            & $g+V+d=L+c+w$ 残差 $<10^{-9}$ & 通过 \\
效率递推            & $E_t-E_{t-1}-0.9c_t+d_t/0.9=0$ & 通过 \\
SOC 边界            & $1200\le E_t\le10800$ & 通过 \\
充放电功率          & $0\le c_t,d_t\le 833.33\kw$ & 通过 \\
充放电互斥          & $c_t\,d_t=0$（逐段检查）& 通过 \\
跨日连续性          & 相邻日末初电量一致（偏差 $<10^{-9}$）& 通过 \\
初末条件            & 起始/终止电量符合设定 & 通过 \\
实际缺口            & $e_t=\max(N_t-z_t,0)$ & 通过 \\
费用复算            & 独立重算与程序输出一致 & 通过 \\
Excel 反读          & 读回结果文件与内部数组一致 & 通过 \\
非前视（信息扰动）  & 扰动未来数据不改变当前决策 & 通过 \\
\bottomrule
\end{tabular}
\end{table}

\textbf{信息扰动测试的具体做法}（以问题三为例）：
在 2 月 1 日、6 月 21 日、12 月 21 日的四个发布时刻，
\textbf{修改当天及未来的全部实际值与未发布预报}，重建预测、残差、融合权重、场景树与 LP，
检查\textbf{当前的预测结果与控制量是否保持不变}。结果显示决策完全不变，
证明程序未使用任何未来信息。

% 按 data/C题.pdf 第1—3页要求，由各问主结果逐段数据重建。
% 仅作汇总，不调用求解器；六个指定十分钟时段取 slot=60,72,84,96,108,120。
\clearpage
\section{题目指定时段与指定日期结果}
\label{sec:required-results}
本附录按题目表1、表2、表3的字段列示五组主结果。电量单位为kWh，费用单位为元；
十分钟时段按起止时刻定位，四小时充放电由连续24段累计，紧急购电按连续发生区间汇总。
第三问与4-3同时列出0点原计划和最终调整后购电，实际总费用包括计划费、净调整费及紧急费。
下列金额和电量由未舍入数据汇总后保留两位小数，舍入后的分项相加可能相差0.01。

\subsection{问题一}
\begin{table}[H]
\centering\footnotesize
\setlength{\tabcolsep}{4pt}
\caption{问题一：指定时段购电及全天汇总（题表1，电量单位：kWh）}
\label{tab:required-q1-grid}
\begin{tabular}{@{}lrlrlr@{}}
\toprule
时间段 & 购电量 & 时间段 & 购电量 & 时间段 & 购电量 \\
\midrule
10:00--10:10 & 0.00 & 12:00--12:10 & 480.41 & 14:00--14:10 & 0.00 \\
16:00--16:10 & 445.43 & 18:00--18:10 & 531.89 & 20:00--20:10 & 0.00 \\
\multicolumn{2}{l}{全天购电量} & 59482.70 & \multicolumn{2}{l}{全天购电费（元）} & 35126.95 \\
\bottomrule
\end{tabular}
\end{table}

\begin{table}[H]
\centering\footnotesize
\setlength{\tabcolsep}{4pt}
\caption{问题一：指定时段充放电及首末储电量（题表2，单位：kWh）}
\label{tab:required-q1-storage}
\begin{tabular}{@{}lrrlrr@{}}
\toprule
时间段 & 充电量 & 放电量 & 时间段 & 充电量 & 放电量 \\
\midrule
00:00--04:00 & 4500.00 & 0.00 & 04:00--08:00 & 833.33 & 6365.84 \\
08:00--12:00 & 4787.96 & 1703.00 & 12:00--16:00 & 5286.04 & 91.10 \\
16:00--20:00 & 0.00 & 5780.13 & 20:00--24:00 & 5333.33 & 2859.87 \\
\multicolumn{2}{l}{0:00储电量} & 6000.00 & \multicolumn{2}{l}{24:00储电量} & 6000.00 \\
\bottomrule
\end{tabular}
\end{table}

\clearpage
\subsection{问题二}
\begin{table}[H]
\centering\footnotesize
\setlength{\tabcolsep}{4pt}
\caption{问题二 2025-03-20：指定时段购电及全天汇总（题表1，电量单位：kWh）}
\label{tab:required-q2-2025-03-20-grid}
\begin{tabular}{@{}lrlrlr@{}}
\toprule
时间段 & 购电量 & 时间段 & 购电量 & 时间段 & 购电量 \\
\midrule
10:00--10:10 & 0.00 & 12:00--12:10 & 626.63 & 14:00--14:10 & 0.00 \\
16:00--16:10 & 494.24 & 18:00--18:10 & 706.68 & 20:00--20:10 & 0.00 \\
\multicolumn{2}{l}{全天计划购电量} & 65615.94 & \multicolumn{2}{l}{计划购电费（元）} & 40493.15 \\
\multicolumn{2}{l}{全天紧急购电量} & 689.80 & \multicolumn{2}{l}{紧急购电费（元）} & 2775.85 \\
\multicolumn{5}{l}{全天总购电费（元）} & 43269.00 \\
\bottomrule
\end{tabular}
\end{table}

\begin{table}[H]
\centering\footnotesize
\setlength{\tabcolsep}{4pt}
\caption{问题二 2025-03-20：指定时段充放电及首末储电量（题表2，单位：kWh）}
\label{tab:required-q2-2025-03-20-storage}
\begin{tabular}{@{}lrrlrr@{}}
\toprule
时间段 & 充电量 & 放电量 & 时间段 & 充电量 & 放电量 \\
\midrule
00:00--04:00 & 1666.67 & 0.00 & 04:00--08:00 & 833.33 & 5501.23 \\
08:00--12:00 & 5912.26 & 3138.77 & 12:00--16:00 & 5454.26 & 566.88 \\
16:00--20:00 & 0.00 & 5671.59 & 20:00--24:00 & 3498.06 & 2968.41 \\
\multicolumn{2}{l}{0:00储电量} & 8550.00 & \multicolumn{2}{l}{24:00储电量} & 4348.26 \\
\bottomrule
\end{tabular}
\end{table}

\FloatBarrier
\begin{table}[H]
\centering\footnotesize
\setlength{\tabcolsep}{4pt}
\caption{问题二 2025-06-21：指定时段购电及全天汇总（题表1，电量单位：kWh）}
\label{tab:required-q2-2025-06-21-grid}
\begin{tabular}{@{}lrlrlr@{}}
\toprule
时间段 & 购电量 & 时间段 & 购电量 & 时间段 & 购电量 \\
\midrule
10:00--10:10 & 0.00 & 12:00--12:10 & 0.00 & 14:00--14:10 & 0.00 \\
16:00--16:10 & 250.70 & 18:00--18:10 & 0.00 & 20:00--20:10 & 0.00 \\
\multicolumn{2}{l}{全天计划购电量} & 46473.41 & \multicolumn{2}{l}{计划购电费（元）} & 25454.37 \\
\multicolumn{2}{l}{全天紧急购电量} & 338.71 & \multicolumn{2}{l}{紧急购电费（元）} & 1155.49 \\
\multicolumn{5}{l}{全天总购电费（元）} & 26609.86 \\
\bottomrule
\end{tabular}
\end{table}

\begin{table}[H]
\centering\footnotesize
\setlength{\tabcolsep}{4pt}
\caption{问题二 2025-06-21：指定时段充放电及首末储电量（题表2，单位：kWh）}
\label{tab:required-q2-2025-06-21-storage}
\begin{tabular}{@{}lrrlrr@{}}
\toprule
时间段 & 充电量 & 放电量 & 时间段 & 充电量 & 放电量 \\
\midrule
00:00--04:00 & 2869.46 & 0.00 & 04:00--08:00 & 833.33 & 2615.32 \\
08:00--12:00 & 5087.51 & 0.00 & 12:00--16:00 & 3470.92 & 0.00 \\
16:00--20:00 & 0.00 & 6017.99 & 20:00--24:00 & 8166.67 & 2622.01 \\
\multicolumn{2}{l}{0:00储电量} & 2670.80 & \multicolumn{2}{l}{24:00储电量} & 8550.00 \\
\bottomrule
\end{tabular}
\end{table}

\FloatBarrier
\begin{table}[H]
\centering\footnotesize
\setlength{\tabcolsep}{4pt}
\caption{问题二 2025-09-23：指定时段购电及全天汇总（题表1，电量单位：kWh）}
\label{tab:required-q2-2025-09-23-grid}
\begin{tabular}{@{}lrlrlr@{}}
\toprule
时间段 & 购电量 & 时间段 & 购电量 & 时间段 & 购电量 \\
\midrule
10:00--10:10 & 0.00 & 12:00--12:10 & 549.48 & 14:00--14:10 & 0.00 \\
16:00--16:10 & 548.33 & 18:00--18:10 & 777.56 & 20:00--20:10 & 0.00 \\
\multicolumn{2}{l}{全天计划购电量} & 70464.22 & \multicolumn{2}{l}{计划购电费（元）} & 43599.64 \\
\multicolumn{2}{l}{全天紧急购电量} & 1179.52 & \multicolumn{2}{l}{紧急购电费（元）} & 5122.78 \\
\multicolumn{5}{l}{全天总购电费（元）} & 48722.41 \\
\bottomrule
\end{tabular}
\end{table}

\begin{table}[H]
\centering\footnotesize
\setlength{\tabcolsep}{4pt}
\caption{问题二 2025-09-23：指定时段充放电及首末储电量（题表2，单位：kWh）}
\label{tab:required-q2-2025-09-23-storage}
\begin{tabular}{@{}lrrlrr@{}}
\toprule
时间段 & 充电量 & 放电量 & 时间段 & 充电量 & 放电量 \\
\midrule
00:00--04:00 & 1666.67 & 0.00 & 04:00--08:00 & 833.33 & 5644.61 \\
08:00--12:00 & 6191.30 & 2995.39 & 12:00--16:00 & 4841.85 & 296.85 \\
16:00--20:00 & 0.00 & 5695.38 & 20:00--24:00 & 8166.67 & 2944.62 \\
\multicolumn{2}{l}{0:00储电量} & 8550.00 & \multicolumn{2}{l}{24:00储电量} & 8550.00 \\
\bottomrule
\end{tabular}
\end{table}

\FloatBarrier
\begin{table}[H]
\centering\footnotesize
\setlength{\tabcolsep}{4pt}
\caption{问题二 2025-12-21：指定时段购电及全天汇总（题表1，电量单位：kWh）}
\label{tab:required-q2-2025-12-21-grid}
\begin{tabular}{@{}lrlrlr@{}}
\toprule
时间段 & 购电量 & 时间段 & 购电量 & 时间段 & 购电量 \\
\midrule
10:00--10:10 & 0.00 & 12:00--12:10 & 989.93 & 14:00--14:10 & 0.00 \\
16:00--16:10 & 863.82 & 18:00--18:10 & 719.77 & 20:00--20:10 & 0.00 \\
\multicolumn{2}{l}{全天计划购电量} & 98935.42 & \multicolumn{2}{l}{计划购电费（元）} & 63572.16 \\
\multicolumn{2}{l}{全天紧急购电量} & 559.54 & \multicolumn{2}{l}{紧急购电费（元）} & 1933.80 \\
\multicolumn{5}{l}{全天总购电费（元）} & 65505.96 \\
\bottomrule
\end{tabular}
\end{table}

\begin{table}[H]
\centering\footnotesize
\setlength{\tabcolsep}{4pt}
\caption{问题二 2025-12-21：指定时段充放电及首末储电量（题表2，单位：kWh）}
\label{tab:required-q2-2025-12-21-storage}
\begin{tabular}{@{}lrrlrr@{}}
\toprule
时间段 & 充电量 & 放电量 & 时间段 & 充电量 & 放电量 \\
\midrule
00:00--04:00 & 1666.67 & 0.00 & 04:00--08:00 & 1666.67 & 1508.33 \\
08:00--12:00 & 5833.33 & 7806.67 & 12:00--16:00 & 8241.15 & 2760.33 \\
16:00--20:00 & 0.00 & 5609.57 & 20:00--24:00 & 8166.67 & 3030.43 \\
\multicolumn{2}{l}{0:00储电量} & 8550.00 & \multicolumn{2}{l}{24:00储电量} & 8550.00 \\
\bottomrule
\end{tabular}
\end{table}

\FloatBarrier
\begin{table}[H]
\centering\footnotesize
\setlength{\tabcolsep}{2.5pt}
\caption{问题二：指定日期全部紧急购电区间（题表3，单位：kWh）}
\label{tab:required-q2-emergency}
\begin{tabular}{@{}lrlrlrlr@{}}
\toprule
\multicolumn{2}{c}{2025.03.20} & \multicolumn{2}{c}{2025.06.21} & \multicolumn{2}{c}{2025.09.23} & \multicolumn{2}{c}{2025.12.21} \\
\midrule
时间段 & 购电量 & 时间段 & 购电量 & 时间段 & 购电量 & 时间段 & 购电量 \\
\midrule
00:20--00:30 & 12.39 & 00:20--00:30 & 40.44 & 00:30--01:20 & 22.17 & 01:00--01:20 & 29.05 \\
01:10--01:40 & 37.48 & 01:20--01:30 & 21.10 & 02:00--02:20 & 55.41 & 02:10--02:20 & 0.61 \\
03:00--03:30 & 70.17 & 01:50--02:10 & 13.24 & 03:50--04:00 & 11.95 & 03:00--03:10 & 15.38 \\
04:50--05:10 & 48.10 & 02:50--03:30 & 99.52 & 04:50--05:10 & 34.28 & 03:30--03:50 & 108.42 \\
05:50--06:00 & 7.85 & 18:50--19:10 & 78.17 & 05:40--05:50 & 10.72 & 05:50--06:00 & 8.04 \\
08:50--09:10 & 26.14 & 20:20--20:30 & 15.05 & 06:20--06:30 & 1.78 & 06:10--06:30 & 41.11 \\
11:10--11:20 & 8.27 & 22:00--22:30 & 71.19 & 10:00--10:20 & 13.05 & 10:30--11:30 & 271.92 \\
12:30--12:50 & 26.69 & — & — & 10:50--11:00 & 5.98 & 15:20--15:30 & 13.09 \\
15:30--16:20 & 208.66 & — & — & 14:30--15:20 & 272.38 & 18:40--19:00 & 8.15 \\
16:30--17:20 & 61.84 & — & — & 15:50--16:50 & 147.40 & 19:40--20:10 & 63.76 \\
18:40--19:10 & 73.67 & — & — & 17:30--17:40 & 16.92 & — & — \\
20:00--20:20 & 41.73 & — & — & 18:10--18:50 & 111.43 & — & — \\
21:10--21:30 & 42.68 & — & — & 19:10--19:30 & 45.49 & — & — \\
23:30--23:50 & 24.14 & — & — & 19:50--20:30 & 140.64 & — & — \\
— & — & — & — & 20:50--21:10 & 106.59 & — & — \\
— & — & — & — & 21:20--22:00 & 58.24 & — & — \\
— & — & — & — & 22:20--22:30 & 57.72 & — & — \\
— & — & — & — & 23:10--23:40 & 67.37 & — & — \\
合计 & 689.80 & 合计 & 338.71 & 合计 & 1179.52 & 合计 & 559.54 \\
\bottomrule
\end{tabular}
\end{table}

\FloatBarrier
\clearpage
\subsection{问题三}
\begin{table}[H]
\centering\footnotesize
\setlength{\tabcolsep}{4pt}
\caption{问题三指定日期费用与紧急电量汇总}
\label{tab:q3-target}
\begin{tabular}{@{}lrrrr@{}}
\toprule
日期 & 原计划费（元） & 净调整费（元） & 紧急电量（kWh） & 总费用（元） \\
\midrule
2025-03-20 & 39999.75 & -94.14 & 873.52 & 43286.36 \\
2025-06-21 & 24123.60 & -105.75 & 342.74 & 25202.45 \\
2025-09-23 & 43523.27 & -155.36 & 949.18 & 47566.49 \\
2025-12-21 & 63209.93 & -75.36 & 607.81 & 65090.39 \\
\bottomrule
\end{tabular}
\end{table}

四个指定日期的净调整费\textbf{均为负}，即调整机制在指定日期上均实现了“省违约金”的净收益。
指定日的完整调度见前文图 \ref{fig:q3-adjust}。

\begin{table}[H]
\centering\footnotesize
\setlength{\tabcolsep}{4pt}
\caption{问题三 2025-03-20：指定时段购电及全天汇总（题表1，电量单位：kWh）}
\label{tab:required-q3-2025-03-20-grid}
\begin{tabular}{@{}lrlrlr@{}}
\toprule
时间段 & 购电量 & 时间段 & 购电量 & 时间段 & 购电量 \\
\midrule
\multicolumn{6}{c}{0点原计划购电量} \\
10:00--10:10 & 0.00 & 12:00--12:10 & 551.93 & 14:00--14:10 & 0.00 \\
16:00--16:10 & 544.34 & 18:00--18:10 & 705.08 & 20:00--20:10 & 0.00 \\
\multicolumn{2}{l}{全天计划购电量} & 64668.95 & \multicolumn{2}{l}{计划购电费（元）} & 39999.75 \\
\multicolumn{6}{c}{最终调整后购电量（不含紧急购电）} \\
10:00--10:10 & 0.00 & 12:00--12:10 & 551.93 & 14:00--14:10 & 0.00 \\
16:00--16:10 & 544.34 & 18:00--18:10 & 705.08 & 20:00--20:10 & 0.00 \\
\multicolumn{2}{l}{全天调整后购电量} & 64101.47 & \multicolumn{2}{l}{净调整费（元）} & -94.14 \\
\multicolumn{2}{l}{全天紧急购电量} & 873.52 & \multicolumn{2}{l}{紧急购电费（元）} & 3380.75 \\
\multicolumn{5}{l}{全天总购电费（元）} & 43286.36 \\
\bottomrule
\end{tabular}
\end{table}

\begin{table}[H]
\centering\footnotesize
\setlength{\tabcolsep}{4pt}
\caption{问题三 2025-03-20：指定时段充放电及首末储电量（题表2，单位：kWh）}
\label{tab:required-q3-2025-03-20-storage}
\begin{tabular}{@{}lrrlrr@{}}
\toprule
时间段 & 充电量 & 放电量 & 时间段 & 充电量 & 放电量 \\
\midrule
00:00--04:00 & 1629.40 & 0.00 & 04:00--08:00 & 833.33 & 5973.45 \\
08:00--12:00 & 6049.15 & 2666.55 & 12:00--16:00 & 5195.34 & 468.04 \\
16:00--20:00 & 2.26 & 5675.55 & 20:00--24:00 & 3395.02 & 2974.39 \\
\multicolumn{2}{l}{0:00储电量} & 8583.54 & \multicolumn{2}{l}{24:00储电量} & 4246.51 \\
\bottomrule
\end{tabular}
\end{table}

\FloatBarrier
\begin{table}[H]
\centering\footnotesize
\setlength{\tabcolsep}{4pt}
\caption{问题三 2025-06-21：指定时段购电及全天汇总（题表1，电量单位：kWh）}
\label{tab:required-q3-2025-06-21-grid}
\begin{tabular}{@{}lrlrlr@{}}
\toprule
时间段 & 购电量 & 时间段 & 购电量 & 时间段 & 购电量 \\
\midrule
\multicolumn{6}{c}{0点原计划购电量} \\
10:00--10:10 & 0.00 & 12:00--12:10 & 0.00 & 14:00--14:10 & 0.00 \\
16:00--16:10 & 223.91 & 18:00--18:10 & 0.00 & 20:00--20:10 & 0.00 \\
\multicolumn{2}{l}{全天计划购电量} & 43570.65 & \multicolumn{2}{l}{计划购电费（元）} & 24123.60 \\
\multicolumn{6}{c}{最终调整后购电量（不含紧急购电）} \\
10:00--10:10 & 0.00 & 12:00--12:10 & 0.00 & 14:00--14:10 & 0.00 \\
16:00--16:10 & 176.27 & 18:00--18:10 & 0.00 & 20:00--20:10 & 0.00 \\
\multicolumn{2}{l}{全天调整后购电量} & 43321.19 & \multicolumn{2}{l}{净调整费（元）} & -105.75 \\
\multicolumn{2}{l}{全天紧急购电量} & 342.74 & \multicolumn{2}{l}{紧急购电费（元）} & 1184.60 \\
\multicolumn{5}{l}{全天总购电费（元）} & 25202.45 \\
\bottomrule
\end{tabular}
\end{table}

\begin{table}[H]
\centering\footnotesize
\setlength{\tabcolsep}{4pt}
\caption{问题三 2025-06-21：指定时段充放电及首末储电量（题表2，单位：kWh）}
\label{tab:required-q3-2025-06-21-storage}
\begin{tabular}{@{}lrrlrr@{}}
\toprule
时间段 & 充电量 & 放电量 & 时间段 & 充电量 & 放电量 \\
\midrule
00:00--04:00 & 762.83 & 0.00 & 04:00--08:00 & 857.88 & 2271.65 \\
08:00--12:00 & 6553.33 & 0.00 & 12:00--16:00 & 4083.07 & 0.00 \\
16:00--20:00 & 0.00 & 5995.04 & 20:00--24:00 & 8149.49 & 2645.72 \\
\multicolumn{2}{l}{0:00储电量} & 2292.66 & \multicolumn{2}{l}{24:00储电量} & 8533.70 \\
\bottomrule
\end{tabular}
\end{table}

\FloatBarrier
\begin{table}[H]
\centering\footnotesize
\setlength{\tabcolsep}{4pt}
\caption{问题三 2025-09-23：指定时段购电及全天汇总（题表1，电量单位：kWh）}
\label{tab:required-q3-2025-09-23-grid}
\begin{tabular}{@{}lrlrlr@{}}
\toprule
时间段 & 购电量 & 时间段 & 购电量 & 时间段 & 购电量 \\
\midrule
\multicolumn{6}{c}{0点原计划购电量} \\
10:00--10:10 & 0.00 & 12:00--12:10 & 502.56 & 14:00--14:10 & 0.00 \\
16:00--16:10 & 560.29 & 18:00--18:10 & 767.73 & 20:00--20:10 & 0.00 \\
\multicolumn{2}{l}{全天计划购电量} & 70219.95 & \multicolumn{2}{l}{计划购电费（元）} & 43523.27 \\
\multicolumn{6}{c}{最终调整后购电量（不含紧急购电）} \\
10:00--10:10 & 0.00 & 12:00--12:10 & 430.27 & 14:00--14:10 & 0.00 \\
16:00--16:10 & 560.29 & 18:00--18:10 & 767.73 & 20:00--20:10 & 0.00 \\
\multicolumn{2}{l}{全天调整后购电量} & 69537.44 & \multicolumn{2}{l}{净调整费（元）} & -155.36 \\
\multicolumn{2}{l}{全天紧急购电量} & 949.18 & \multicolumn{2}{l}{紧急购电费（元）} & 4198.58 \\
\multicolumn{5}{l}{全天总购电费（元）} & 47566.49 \\
\bottomrule
\end{tabular}
\end{table}

\begin{table}[H]
\centering\footnotesize
\setlength{\tabcolsep}{4pt}
\caption{问题三 2025-09-23：指定时段充放电及首末储电量（题表2，单位：kWh）}
\label{tab:required-q3-2025-09-23-storage}
\begin{tabular}{@{}lrrlrr@{}}
\toprule
时间段 & 充电量 & 放电量 & 时间段 & 充电量 & 放电量 \\
\midrule
00:00--04:00 & 1682.24 & 0.00 & 04:00--08:00 & 833.33 & 5844.03 \\
08:00--12:00 & 6049.67 & 2598.82 & 12:00--16:00 & 4754.75 & 331.14 \\
16:00--20:00 & 42.16 & 5706.45 & 20:00--24:00 & 8152.71 & 2962.87 \\
\multicolumn{2}{l}{0:00储电量} & 8535.99 & \multicolumn{2}{l}{24:00储电量} & 8517.90 \\
\bottomrule
\end{tabular}
\end{table}

\FloatBarrier
\begin{table}[H]
\centering\footnotesize
\setlength{\tabcolsep}{4pt}
\caption{问题三 2025-12-21：指定时段购电及全天汇总（题表1，电量单位：kWh）}
\label{tab:required-q3-2025-12-21-grid}
\begin{tabular}{@{}lrlrlr@{}}
\toprule
时间段 & 购电量 & 时间段 & 购电量 & 时间段 & 购电量 \\
\midrule
\multicolumn{6}{c}{0点原计划购电量} \\
10:00--10:10 & 0.00 & 12:00--12:10 & 965.84 & 14:00--14:10 & 0.00 \\
16:00--16:10 & 872.53 & 18:00--18:10 & 708.71 & 20:00--20:10 & 0.00 \\
\multicolumn{2}{l}{全天计划购电量} & 98380.06 & \multicolumn{2}{l}{计划购电费（元）} & 63209.93 \\
\multicolumn{6}{c}{最终调整后购电量（不含紧急购电）} \\
10:00--10:10 & 0.00 & 12:00--12:10 & 954.33 & 14:00--14:10 & 0.00 \\
16:00--16:10 & 872.53 & 18:00--18:10 & 708.71 & 20:00--20:10 & 0.00 \\
\multicolumn{2}{l}{全天调整后购电量} & 98056.59 & \multicolumn{2}{l}{净调整费（元）} & -75.36 \\
\multicolumn{2}{l}{全天紧急购电量} & 607.81 & \multicolumn{2}{l}{紧急购电费（元）} & 1955.82 \\
\multicolumn{5}{l}{全天总购电费（元）} & 65090.39 \\
\bottomrule
\end{tabular}
\end{table}

\begin{table}[H]
\centering\footnotesize
\setlength{\tabcolsep}{4pt}
\caption{问题三 2025-12-21：指定时段充放电及首末储电量（题表2，单位：kWh）}
\label{tab:required-q3-2025-12-21-storage}
\begin{tabular}{@{}lrrlrr@{}}
\toprule
时间段 & 充电量 & 放电量 & 时间段 & 充电量 & 放电量 \\
\midrule
00:00--04:00 & 1640.14 & 0.00 & 04:00--08:00 & 1660.82 & 1605.72 \\
08:00--12:00 & 5884.07 & 7358.10 & 12:00--16:00 & 7766.93 & 2767.14 \\
16:00--20:00 & 7.42 & 5602.70 & 20:00--24:00 & 8200.89 & 3043.38 \\
\multicolumn{2}{l}{0:00储电量} & 8573.87 & \multicolumn{2}{l}{24:00储电量} & 8576.96 \\
\bottomrule
\end{tabular}
\end{table}

\FloatBarrier
\begin{table}[H]
\centering\footnotesize
\setlength{\tabcolsep}{2.5pt}
\caption{问题三：指定日期全部紧急购电区间（题表3，单位：kWh）}
\label{tab:required-q3-emergency}
\begin{tabular}{@{}lrlrlrlr@{}}
\toprule
\multicolumn{2}{c}{2025.03.20} & \multicolumn{2}{c}{2025.06.21} & \multicolumn{2}{c}{2025.09.23} & \multicolumn{2}{c}{2025.12.21} \\
\midrule
时间段 & 购电量 & 时间段 & 购电量 & 时间段 & 购电量 & 时间段 & 购电量 \\
\midrule
00:20--00:30 & 12.39 & 00:20--00:30 & 40.44 & 00:30--01:20 & 22.17 & 01:00--01:20 & 29.05 \\
01:10--01:40 & 37.48 & 01:20--01:30 & 21.10 & 02:00--02:20 & 55.41 & 02:10--02:20 & 0.61 \\
03:00--03:30 & 70.17 & 01:50--02:10 & 13.24 & 03:50--04:00 & 11.95 & 03:00--03:10 & 15.38 \\
04:50--05:10 & 48.10 & 02:50--03:30 & 99.52 & 04:50--05:10 & 34.31 & 03:30--03:50 & 108.42 \\
05:50--06:00 & 6.72 & 18:50--19:10 & 79.29 & 05:40--05:50 & 17.70 & 05:50--06:00 & 8.04 \\
08:10--09:20 & 147.65 & 20:20--20:30 & 18.51 & 06:10--06:30 & 0.57 & 06:20--06:30 & 15.63 \\
09:40--10:00 & 56.01 & 22:00--22:30 & 70.63 & 07:40--07:50 & 5.03 & 10:30--11:40 & 363.36 \\
11:00--11:30 & 114.30 & — & — & 14:20--15:00 & 75.47 & 13:00--13:10 & 1.78 \\
11:40--12:00 & 35.27 & — & — & 15:50--16:40 & 126.89 & 18:40--19:00 & 6.49 \\
12:30--12:50 & 109.16 & — & — & 17:30--17:50 & 19.82 & 19:40--20:00 & 59.05 \\
13:10--13:20 & 0.81 & — & — & 18:10--18:50 & 96.87 & — & — \\
14:00--14:10 & 6.50 & — & — & 19:10--19:30 & 41.82 & — & — \\
15:50--16:10 & 22.60 & — & — & 19:50--20:30 & 131.14 & — & — \\
16:50--17:00 & 0.11 & — & — & 20:50--21:10 & 106.59 & — & — \\
18:40--19:10 & 85.04 & — & — & 21:20--22:00 & 78.37 & — & — \\
19:50--20:20 & 49.50 & — & — & 22:20--22:30 & 57.72 & — & — \\
21:10--21:30 & 47.58 & — & — & 23:10--23:40 & 67.37 & — & — \\
23:30--23:50 & 24.14 & — & — & — & — & — & — \\
合计 & 873.52 & 合计 & 342.74 & 合计 & 949.18 & 合计 & 607.81 \\
\bottomrule
\end{tabular}
\end{table}

\FloatBarrier
\clearpage
\subsection{问题四（4-2）}
\begin{table}[H]
\centering\footnotesize
\setlength{\tabcolsep}{4pt}
\caption{问题四（4-2） 2025-03-20：指定时段购电及全天汇总（题表1，电量单位：kWh）}
\label{tab:required-q42-2025-03-20-grid}
\begin{tabular}{@{}lrlrlr@{}}
\toprule
时间段 & 购电量 & 时间段 & 购电量 & 时间段 & 购电量 \\
\midrule
10:00--10:10 & 0.00 & 12:00--12:10 & 653.15 & 14:00--14:10 & 12.32 \\
16:00--16:10 & 494.35 & 18:00--18:10 & 0.00 & 20:00--20:10 & 731.74 \\
\multicolumn{2}{l}{全天计划购电量} & 63984.26 & \multicolumn{2}{l}{计划购电费（元）} & 41035.65 \\
\multicolumn{2}{l}{全天紧急购电量} & 631.45 & \multicolumn{2}{l}{紧急购电费（元）} & 2706.13 \\
\multicolumn{5}{l}{全天总购电费（元）} & 43741.78 \\
\bottomrule
\end{tabular}
\end{table}

\begin{table}[H]
\centering\footnotesize
\setlength{\tabcolsep}{4pt}
\caption{问题四（4-2） 2025-03-20：指定时段充放电及首末储电量（题表2，单位：kWh）}
\label{tab:required-q42-2025-03-20-storage}
\begin{tabular}{@{}lrrlrr@{}}
\toprule
时间段 & 充电量 & 放电量 & 时间段 & 充电量 & 放电量 \\
\midrule
00:00--04:00 & 2500.00 & 0.00 & 04:00--08:00 & 1500.00 & 5904.30 \\
08:00--12:00 & 5922.21 & 3144.53 & 12:00--16:00 & 5058.38 & 520.46 \\
16:00--20:00 & 0.00 & 6405.44 & 20:00--24:00 & 833.33 & 2234.56 \\
\multicolumn{2}{l}{0:00储电量} & 7950.00 & \multicolumn{2}{l}{24:00储电量} & 1950.00 \\
\bottomrule
\end{tabular}
\end{table}

\FloatBarrier
\begin{table}[H]
\centering\footnotesize
\setlength{\tabcolsep}{4pt}
\caption{问题四（4-2） 2025-06-21：指定时段购电及全天汇总（题表1，电量单位：kWh）}
\label{tab:required-q42-2025-06-21-grid}
\begin{tabular}{@{}lrlrlr@{}}
\toprule
时间段 & 购电量 & 时间段 & 购电量 & 时间段 & 购电量 \\
\midrule
10:00--10:10 & 0.00 & 12:00--12:10 & 0.00 & 14:00--14:10 & 0.00 \\
16:00--16:10 & 252.87 & 18:00--18:10 & 487.95 & 20:00--20:10 & 0.00 \\
\multicolumn{2}{l}{全天计划购电量} & 49515.72 & \multicolumn{2}{l}{计划购电费（元）} & 23064.36 \\
\multicolumn{2}{l}{全天紧急购电量} & 311.29 & \multicolumn{2}{l}{紧急购电费（元）} & 1037.19 \\
\multicolumn{5}{l}{全天总购电费（元）} & 24101.55 \\
\bottomrule
\end{tabular}
\end{table}

\begin{table}[H]
\centering\footnotesize
\setlength{\tabcolsep}{4pt}
\caption{问题四（4-2） 2025-06-21：指定时段充放电及首末储电量（题表2，单位：kWh）}
\label{tab:required-q42-2025-06-21-storage}
\begin{tabular}{@{}lrrlrr@{}}
\toprule
时间段 & 充电量 & 放电量 & 时间段 & 充电量 & 放电量 \\
\midrule
00:00--04:00 & 0.00 & 0.00 & 04:00--08:00 & 2749.81 & 2227.34 \\
08:00--12:00 & 7535.00 & 0.00 & 12:00--16:00 & 3131.66 & 0.00 \\
16:00--20:00 & 0.00 & 6017.99 & 20:00--24:00 & 9166.67 & 2757.01 \\
\multicolumn{2}{l}{0:00储电量} & 1200.00 & \multicolumn{2}{l}{24:00储电量} & 9300.00 \\
\bottomrule
\end{tabular}
\end{table}

\FloatBarrier
\begin{table}[H]
\centering\footnotesize
\setlength{\tabcolsep}{4pt}
\caption{问题四（4-2） 2025-09-23：指定时段购电及全天汇总（题表1，电量单位：kWh）}
\label{tab:required-q42-2025-09-23-grid}
\begin{tabular}{@{}lrlrlr@{}}
\toprule
时间段 & 购电量 & 时间段 & 购电量 & 时间段 & 购电量 \\
\midrule
10:00--10:10 & 0.00 & 12:00--12:10 & 568.09 & 14:00--14:10 & 0.00 \\
16:00--16:10 & 550.06 & 18:00--18:10 & 0.00 & 20:00--20:10 & 0.00 \\
\multicolumn{2}{l}{全天计划购电量} & 70272.43 & \multicolumn{2}{l}{计划购电费（元）} & 45565.16 \\
\multicolumn{2}{l}{全天紧急购电量} & 1064.22 & \multicolumn{2}{l}{紧急购电费（元）} & 4972.92 \\
\multicolumn{5}{l}{全天总购电费（元）} & 50538.08 \\
\bottomrule
\end{tabular}
\end{table}

\begin{table}[H]
\centering\footnotesize
\setlength{\tabcolsep}{4pt}
\caption{问题四（4-2） 2025-09-23：指定时段充放电及首末储电量（题表2，单位：kWh）}
\label{tab:required-q42-2025-09-23-storage}
\begin{tabular}{@{}lrrlrr@{}}
\toprule
时间段 & 充电量 & 放电量 & 时间段 & 充电量 & 放电量 \\
\midrule
00:00--04:00 & 4166.67 & 0.00 & 04:00--08:00 & 833.33 & 5614.15 \\
08:00--12:00 & 6150.98 & 3025.85 & 12:00--16:00 & 4882.17 & 296.85 \\
16:00--20:00 & 0.00 & 6245.22 & 20:00--24:00 & 4833.33 & 2394.78 \\
\multicolumn{2}{l}{0:00储电量} & 6300.00 & \multicolumn{2}{l}{24:00储电量} & 5550.00 \\
\bottomrule
\end{tabular}
\end{table}

\FloatBarrier
\begin{table}[H]
\centering\footnotesize
\setlength{\tabcolsep}{4pt}
\caption{问题四（4-2） 2025-12-21：指定时段购电及全天汇总（题表1，电量单位：kWh）}
\label{tab:required-q42-2025-12-21-grid}
\begin{tabular}{@{}lrlrlr@{}}
\toprule
时间段 & 购电量 & 时间段 & 购电量 & 时间段 & 购电量 \\
\midrule
10:00--10:10 & 0.00 & 12:00--12:10 & 993.19 & 14:00--14:10 & 0.00 \\
16:00--16:10 & 863.82 & 18:00--18:10 & 719.77 & 20:00--20:10 & 0.00 \\
\multicolumn{2}{l}{全天计划购电量} & 97387.54 & \multicolumn{2}{l}{计划购电费（元）} & 74053.41 \\
\multicolumn{2}{l}{全天紧急购电量} & 526.64 & \multicolumn{2}{l}{紧急购电费（元）} & 2403.13 \\
\multicolumn{5}{l}{全天总购电费（元）} & 76456.54 \\
\bottomrule
\end{tabular}
\end{table}

\begin{table}[H]
\centering\footnotesize
\setlength{\tabcolsep}{4pt}
\caption{问题四（4-2） 2025-12-21：指定时段充放电及首末储电量（题表2，单位：kWh）}
\label{tab:required-q42-2025-12-21-storage}
\begin{tabular}{@{}lrrlrr@{}}
\toprule
时间段 & 充电量 & 放电量 & 时间段 & 充电量 & 放电量 \\
\midrule
00:00--04:00 & 666.67 & 0.00 & 04:00--08:00 & 0.00 & 833.33 \\
08:00--12:00 & 5833.33 & 7806.67 & 12:00--16:00 & 8256.22 & 2772.54 \\
16:00--20:00 & 0.00 & 5609.30 & 20:00--24:00 & 8333.33 & 3030.70 \\
\multicolumn{2}{l}{0:00储电量} & 10200.00 & \multicolumn{2}{l}{24:00储电量} & 8700.00 \\
\bottomrule
\end{tabular}
\end{table}

\FloatBarrier
\begin{table}[H]
\centering\footnotesize
\setlength{\tabcolsep}{2.5pt}
\caption{问题四（4-2）：指定日期全部紧急购电区间（题表3，单位：kWh）}
\label{tab:required-q42-emergency}
\begin{tabular}{@{}lrlrlrlr@{}}
\toprule
\multicolumn{2}{c}{2025.03.20} & \multicolumn{2}{c}{2025.06.21} & \multicolumn{2}{c}{2025.09.23} & \multicolumn{2}{c}{2025.12.21} \\
\midrule
时间段 & 购电量 & 时间段 & 购电量 & 时间段 & 购电量 & 时间段 & 购电量 \\
\midrule
00:20--00:30 & 10.64 & 00:20--00:30 & 36.51 & 00:30--00:50 & 10.00 & 01:00--01:20 & 26.83 \\
01:10--01:40 & 32.72 & 01:20--01:30 & 16.72 & 01:00--01:10 & 5.36 & 03:00--03:10 & 13.56 \\
03:00--03:30 & 67.32 & 02:00--02:10 & 10.40 & 02:00--02:20 & 52.79 & 03:30--03:50 & 103.82 \\
04:50--05:10 & 47.52 & 02:50--03:30 & 90.63 & 03:50--04:00 & 9.70 & 05:50--06:00 & 4.79 \\
05:50--06:00 & 6.72 & 18:50--19:10 & 78.17 & 04:50--05:10 & 31.67 & 06:10--06:30 & 34.39 \\
09:00--09:10 & 17.38 & 20:20--20:30 & 15.05 & 05:40--05:50 & 7.01 & 10:30--11:30 & 259.80 \\
12:30--12:50 & 3.28 & 22:00--22:30 & 63.81 & 10:00--10:20 & 4.62 & 15:20--15:30 & 11.52 \\
15:30--16:20 & 199.58 & — & — & 14:30--15:20 & 260.90 & 18:40--19:00 & 8.15 \\
16:30--17:20 & 61.32 & — & — & 15:50--16:50 & 122.31 & 19:40--20:10 & 63.76 \\
18:40--19:10 & 78.09 & — & — & 17:30--17:40 & 16.92 & — & — \\
20:00--20:20 & 43.60 & — & — & 18:10--18:50 & 96.79 & — & — \\
21:10--21:30 & 40.07 & — & — & 19:10--19:30 & 41.85 & — & — \\
23:30--23:50 & 23.20 & — & — & 19:50--20:30 & 140.17 & — & — \\
— & — & — & — & 20:50--21:10 & 105.01 & — & — \\
— & — & — & — & 21:20--22:00 & 54.90 & — & — \\
— & — & — & — & 22:20--22:30 & 57.72 & — & — \\
— & — & — & — & 23:10--23:40 & 46.50 & — & — \\
合计 & 631.45 & 合计 & 311.29 & 合计 & 1064.22 & 合计 & 526.64 \\
\bottomrule
\end{tabular}
\end{table}

\FloatBarrier
\clearpage
\subsection{问题四（4-3）}
\begin{table}[H]
\centering\footnotesize
\setlength{\tabcolsep}{4pt}
\caption{问题四（4-3） 2025-03-20：指定时段购电及全天汇总（题表1，电量单位：kWh）}
\label{tab:required-q43-2025-03-20-grid}
\begin{tabular}{@{}lrlrlr@{}}
\toprule
时间段 & 购电量 & 时间段 & 购电量 & 时间段 & 购电量 \\
\midrule
\multicolumn{6}{c}{0点原计划购电量} \\
10:00--10:10 & 0.00 & 12:00--12:10 & 545.41 & 14:00--14:10 & 0.00 \\
16:00--16:10 & 525.49 & 18:00--18:10 & 0.00 & 20:00--20:10 & 719.50 \\
\multicolumn{2}{l}{全天计划购电量} & 67104.27 & \multicolumn{2}{l}{计划购电费（元）} & 41826.94 \\
\multicolumn{6}{c}{最终调整后购电量（不含紧急购电）} \\
10:00--10:10 & 0.00 & 12:00--12:10 & 545.41 & 14:00--14:10 & 0.00 \\
16:00--16:10 & 525.49 & 18:00--18:10 & 0.00 & 20:00--20:10 & 719.50 \\
\multicolumn{2}{l}{全天调整后购电量} & 66322.18 & \multicolumn{2}{l}{净调整费（元）} & -164.65 \\
\multicolumn{2}{l}{全天紧急购电量} & 1936.99 & \multicolumn{2}{l}{紧急购电费（元）} & 7681.69 \\
\multicolumn{5}{l}{全天总购电费（元）} & 49343.99 \\
\bottomrule
\end{tabular}
\end{table}

\begin{table}[H]
\centering\footnotesize
\setlength{\tabcolsep}{4pt}
\caption{问题四（4-3） 2025-03-20：指定时段充放电及首末储电量（题表2，单位：kWh）}
\label{tab:required-q43-2025-03-20-storage}
\begin{tabular}{@{}lrrlrr@{}}
\toprule
时间段 & 充电量 & 放电量 & 时间段 & 充电量 & 放电量 \\
\midrule
00:00--04:00 & 7500.00 & 0.00 & 04:00--08:00 & 3166.67 & 6274.00 \\
08:00--12:00 & 5993.55 & 2366.00 & 12:00--16:00 & 5097.01 & 343.95 \\
16:00--20:00 & 0.73 & 7170.36 & 20:00--24:00 & 65.10 & 1522.37 \\
\multicolumn{2}{l}{0:00储电量} & 1200.00 & \multicolumn{2}{l}{24:00储电量} & 1200.00 \\
\bottomrule
\end{tabular}
\end{table}

\FloatBarrier
\begin{table}[H]
\centering\footnotesize
\setlength{\tabcolsep}{4pt}
\caption{问题四（4-3） 2025-06-21：指定时段购电及全天汇总（题表1，电量单位：kWh）}
\label{tab:required-q43-2025-06-21-grid}
\begin{tabular}{@{}lrlrlr@{}}
\toprule
时间段 & 购电量 & 时间段 & 购电量 & 时间段 & 购电量 \\
\midrule
\multicolumn{6}{c}{0点原计划购电量} \\
10:00--10:10 & 0.00 & 12:00--12:10 & 0.00 & 14:00--14:10 & 0.00 \\
16:00--16:10 & 192.84 & 18:00--18:10 & 453.05 & 20:00--20:10 & 0.00 \\
\multicolumn{2}{l}{全天计划购电量} & 34666.48 & \multicolumn{2}{l}{计划购电费（元）} & 17225.37 \\
\multicolumn{6}{c}{最终调整后购电量（不含紧急购电）} \\
10:00--10:10 & 0.00 & 12:00--12:10 & 0.00 & 14:00--14:10 & 0.00 \\
16:00--16:10 & 134.07 & 18:00--18:10 & 453.05 & 20:00--20:10 & 0.00 \\
\multicolumn{2}{l}{全天调整后购电量} & 34353.81 & \multicolumn{2}{l}{净调整费（元）} & -93.39 \\
\multicolumn{2}{l}{全天紧急购电量} & 659.63 & \multicolumn{2}{l}{紧急购电费（元）} & 2033.32 \\
\multicolumn{5}{l}{全天总购电费（元）} & 19165.30 \\
\bottomrule
\end{tabular}
\end{table}

\begin{table}[H]
\centering\footnotesize
\setlength{\tabcolsep}{4pt}
\caption{问题四（4-3） 2025-06-21：指定时段充放电及首末储电量（题表2，单位：kWh）}
\label{tab:required-q43-2025-06-21-storage}
\begin{tabular}{@{}lrrlrr@{}}
\toprule
时间段 & 充电量 & 放电量 & 时间段 & 充电量 & 放电量 \\
\midrule
00:00--04:00 & 0.00 & 0.00 & 04:00--08:00 & 2586.28 & 2075.00 \\
08:00--12:00 & 6553.33 & 0.00 & 12:00--16:00 & 4107.81 & 15.41 \\
16:00--20:00 & 0.00 & 6096.96 & 20:00--24:00 & 53.86 & 2586.67 \\
\multicolumn{2}{l}{0:00储电量} & 1200.00 & \multicolumn{2}{l}{24:00储电量} & 1200.00 \\
\bottomrule
\end{tabular}
\end{table}

\FloatBarrier
\begin{table}[H]
\centering\footnotesize
\setlength{\tabcolsep}{4pt}
\caption{问题四（4-3） 2025-09-23：指定时段购电及全天汇总（题表1，电量单位：kWh）}
\label{tab:required-q43-2025-09-23-grid}
\begin{tabular}{@{}lrlrlr@{}}
\toprule
时间段 & 购电量 & 时间段 & 购电量 & 时间段 & 购电量 \\
\midrule
\multicolumn{6}{c}{0点原计划购电量} \\
10:00--10:10 & 0.00 & 12:00--12:10 & 456.65 & 14:00--14:10 & 0.00 \\
16:00--16:10 & 536.21 & 18:00--18:10 & 0.00 & 20:00--20:10 & 0.00 \\
\multicolumn{2}{l}{全天计划购电量} & 66826.12 & \multicolumn{2}{l}{计划购电费（元）} & 43138.23 \\
\multicolumn{6}{c}{最终调整后购电量（不含紧急购电）} \\
10:00--10:10 & 0.00 & 12:00--12:10 & 306.14 & 14:00--14:10 & 0.00 \\
16:00--16:10 & 536.21 & 18:00--18:10 & 0.00 & 20:00--20:10 & 0.00 \\
\multicolumn{2}{l}{全天调整后购电量} & 66054.20 & \multicolumn{2}{l}{净调整费（元）} & -57.05 \\
\multicolumn{2}{l}{全天紧急购电量} & 2121.28 & \multicolumn{2}{l}{紧急购电费（元）} & 9624.84 \\
\multicolumn{5}{l}{全天总购电费（元）} & 52706.02 \\
\bottomrule
\end{tabular}
\end{table}

\begin{table}[H]
\centering\footnotesize
\setlength{\tabcolsep}{4pt}
\caption{问题四（4-3） 2025-09-23：指定时段充放电及首末储电量（题表2，单位：kWh）}
\label{tab:required-q43-2025-09-23-storage}
\begin{tabular}{@{}lrrlrr@{}}
\toprule
时间段 & 充电量 & 放电量 & 时间段 & 充电量 & 放电量 \\
\midrule
00:00--04:00 & 7333.33 & 0.00 & 04:00--08:00 & 3333.33 & 6441.96 \\
08:00--12:00 & 6025.31 & 2198.04 & 12:00--16:00 & 4936.43 & 245.51 \\
16:00--20:00 & 25.85 & 6307.45 & 20:00--24:00 & 16.57 & 2360.41 \\
\multicolumn{2}{l}{0:00储电量} & 1200.00 & \multicolumn{2}{l}{24:00储电量} & 1200.00 \\
\bottomrule
\end{tabular}
\end{table}

\FloatBarrier
\begin{table}[H]
\centering\footnotesize
\setlength{\tabcolsep}{4pt}
\caption{问题四（4-3） 2025-12-21：指定时段购电及全天汇总（题表1，电量单位：kWh）}
\label{tab:required-q43-2025-12-21-grid}
\begin{tabular}{@{}lrlrlr@{}}
\toprule
时间段 & 购电量 & 时间段 & 购电量 & 时间段 & 购电量 \\
\midrule
\multicolumn{6}{c}{0点原计划购电量} \\
10:00--10:10 & 0.00 & 12:00--12:10 & 918.26 & 14:00--14:10 & 0.00 \\
16:00--16:10 & 843.15 & 18:00--18:10 & 703.41 & 20:00--20:10 & 0.00 \\
\multicolumn{2}{l}{全天计划购电量} & 95873.43 & \multicolumn{2}{l}{计划购电费（元）} & 72035.15 \\
\multicolumn{6}{c}{最终调整后购电量（不含紧急购电）} \\
10:00--10:10 & 0.00 & 12:00--12:10 & 918.26 & 14:00--14:10 & 0.00 \\
16:00--16:10 & 844.01 & 18:00--18:10 & 703.41 & 20:00--20:10 & 0.00 \\
\multicolumn{2}{l}{全天调整后购电量} & 95320.49 & \multicolumn{2}{l}{净调整费（元）} & -205.31 \\
\multicolumn{2}{l}{全天紧急购电量} & 1295.96 & \multicolumn{2}{l}{紧急购电费（元）} & 5831.56 \\
\multicolumn{5}{l}{全天总购电费（元）} & 77661.40 \\
\bottomrule
\end{tabular}
\end{table}

\begin{table}[H]
\centering\footnotesize
\setlength{\tabcolsep}{4pt}
\caption{问题四（4-3） 2025-12-21：指定时段充放电及首末储电量（题表2，单位：kWh）}
\label{tab:required-q43-2025-12-21-storage}
\begin{tabular}{@{}lrrlrr@{}}
\toprule
时间段 & 充电量 & 放电量 & 时间段 & 充电量 & 放电量 \\
\midrule
00:00--04:00 & 10666.67 & 0.00 & 04:00--08:00 & 3.60 & 864.86 \\
08:00--12:00 & 5833.33 & 7351.09 & 12:00--16:00 & 7885.43 & 2899.17 \\
16:00--20:00 & 12.56 & 5655.47 & 20:00--24:00 & 44.62 & 3030.85 \\
\multicolumn{2}{l}{0:00储电量} & 1200.00 & \multicolumn{2}{l}{24:00储电量} & 1200.00 \\
\bottomrule
\end{tabular}
\end{table}

\FloatBarrier
\begin{table}[H]
\centering\footnotesize
\setlength{\tabcolsep}{2.5pt}
\caption{问题四（4-3）：指定日期全部紧急购电区间（题表3，单位：kWh）}
\label{tab:required-q43-emergency}
\begin{tabular}{@{}lrlrlrlr@{}}
\toprule
\multicolumn{2}{c}{2025.03.20} & \multicolumn{2}{c}{2025.06.21} & \multicolumn{2}{c}{2025.09.23} & \multicolumn{2}{c}{2025.12.21} \\
\midrule
时间段 & 购电量 & 时间段 & 购电量 & 时间段 & 购电量 & 时间段 & 购电量 \\
\midrule
00:20--00:30 & 22.46 & 00:10--00:30 & 74.72 & 00:30--01:20 & 111.02 & 00:00--00:20 & 7.92 \\
01:00--02:00 & 82.55 & 01:20--02:20 & 123.00 & 01:40--02:30 & 112.37 & 01:00--01:20 & 57.62 \\
02:30--02:40 & 0.58 & 02:50--03:30 & 149.65 & 03:40--04:00 & 34.84 & 01:30--01:40 & 11.18 \\
03:00--03:40 & 111.77 & 03:50--04:10 & 13.90 & 04:50--05:20 & 67.28 & 02:10--02:20 & 15.66 \\
04:50--05:20 & 81.18 & 17:20--17:40 & 14.78 & 05:30--05:50 & 29.41 & 02:40--02:50 & 0.88 \\
05:50--06:00 & 20.46 & 18:50--19:10 & 105.08 & 06:10--06:30 & 28.01 & 03:00--03:20 & 37.48 \\
06:40--07:00 & 7.10 & 20:20--20:40 & 52.87 & 06:50--07:10 & 12.98 & 03:30--03:50 & 144.79 \\
08:00--09:20 & 342.10 & 21:40--22:40 & 125.64 & 07:40--08:00 & 26.75 & 04:00--04:20 & 11.40 \\
09:40--10:00 & 80.43 & — & — & 09:30--09:50 & 23.79 & 05:40--06:30 & 97.47 \\
10:20--10:30 & 17.89 & — & — & 10:00--10:20 & 7.35 & 07:10--07:20 & 3.74 \\
10:50--12:00 & 461.66 & — & — & 10:40--11:00 & 35.51 & 07:50--08:00 & 2.03 \\
12:30--13:30 & 191.60 & — & — & 11:30--11:50 & 22.71 & 08:30--09:00 & 24.22 \\
14:00--14:20 & 35.44 & — & — & 12:50--13:00 & 0.16 & 09:30--09:50 & 39.29 \\
14:50--15:00 & 7.47 & — & — & 13:50--14:00 & 14.31 & 10:20--11:50 & 574.75 \\
15:30--16:20 & 90.12 & — & — & 14:20--15:20 & 217.15 & 12:00--12:20 & 39.04 \\
16:40--17:00 & 27.47 & — & — & 15:50--16:50 & 263.19 & 12:40--13:10 & 63.72 \\
18:40--19:20 & 127.95 & — & — & 17:00--17:50 & 73.03 & 15:20--15:30 & 5.86 \\
19:30--20:20 & 93.09 & — & — & 18:00--18:50 & 195.57 & 16:20--16:30 & 11.34 \\
21:10--21:40 & 75.12 & — & — & 19:10--19:30 & 82.27 & 18:00--18:20 & 11.27 \\
23:30--23:50 & 60.55 & — & — & 19:50--22:00 & 562.42 & 18:40--19:00 & 30.66 \\
— & — & — & — & 22:10--22:40 & 77.99 & 19:40--20:10 & 88.22 \\
— & — & — & — & 23:10--23:40 & 123.17 & 21:50--22:00 & 2.84 \\
— & — & — & — & — & — & 22:50--23:00 & 12.06 \\
— & — & — & — & — & — & 23:30--23:40 & 2.50 \\
合计 & 1936.99 & 合计 & 659.63 & 合计 & 2121.28 & 合计 & 1295.96 \\
\bottomrule
\end{tabular}
\end{table}

\FloatBarrier


% ==================== 附录：完整可运行源程序 ====================
% 按 format2026 第五条要求，附录包含建模所用到的完整、可运行的源程序。
% 本文件由 scripts/gen_appendix_source.py 自动生成，勿手改。
% 源码直接取自支撑材料文件夹 final_version/，故与所提交文件逐字一致。
\section{建模源程序（完整可运行）}
\label{sec:appendix-source}

以下按问题列出\textbf{生成论文全部数值结果}所需源程序。
这些代码\textbf{直接取自支撑材料文件夹 \texttt{final\_version/}}，
因此与所提交文件\textbf{逐字一致}；
运行环境与复现命令见 \S\ref{sec:appendix-env}。
全部程序使用 Python 语言与 \texttt{numpy}/\texttt{scipy}（HiGHS 求解器），
Excel 读写使用 \texttt{openpyxl}；未使用 SPSS 等需手工交互的软件，
故无交互命令需要单独记录。

\textbf{收录边界}：本文的图表由配图脚本从上述程序保存的结果文件
（\texttt{*.npz}/\texttt{*.csv}）重建，属于表达层，
不参与任何数值计算，也不影响任何结果；
按“仅收录与建模结果直接相关的源程序”的原则，
\texttt{plot\_*.py} 与 \texttt{common/plotting/} 未列入本附录，
其完整文件已随支撑材料一并提交。

\begingroup
\footnotesize
\setlength{\parindent}{0pt}
\setlength{\parskip}{0pt}

\subsection{问题一：确定性调度}

\subsubsection*{q1\_model.py}
\verbatiminput{../final_version/Q1/q1_code/q1_model.py}

\subsubsection*{run\_q1.py}
\verbatiminput{../final_version/Q1/q1_code/run_q1.py}

\subsection{问题二：滚动随机日前调度}

\subsubsection*{rolling\_window.py}
\verbatiminput{../final_version/Q2/code/rolling_window.py}

\subsubsection*{run\_experiment.py}
\verbatiminput{../final_version/Q2/code/run_experiment.py}

\subsubsection*{verify\_window.py}
\verbatiminput{../final_version/Q2/code/verify_window.py}

\subsubsection*{check\_information.py}
\verbatiminput{../final_version/Q2/code/check_information.py}

\subsubsection*{report\_experiment.py}
\verbatiminput{../final_version/Q2/code/report_experiment.py}

\subsubsection*{make\_target\_tables.py}
\verbatiminput{../final_version/Q2/code/make_target_tables.py}

\subsection{问题三：预报融合与场景树多时点调整}

\subsubsection*{q3\_data.py}
\verbatiminput{../final_version/Q3/code/q3_data.py}

\subsubsection*{q3\_forecast.py}
\verbatiminput{../final_version/Q3/code/q3_forecast.py}

\subsubsection*{q3\_model.py}
\verbatiminput{../final_version/Q3/code/q3_model.py}

\subsubsection*{run\_q3.py}
\verbatiminput{../final_version/Q3/code/run_q3.py}

\subsubsection*{verify\_q3.py}
\verbatiminput{../final_version/Q3/code/verify_q3.py}

\subsection{问题四-2：波动电价下的日前调度}

\subsubsection*{q42\_model.py}
\verbatiminput{../final_version/Q4/part2/code/q42_model.py}

\subsubsection*{run\_q42.py}
\verbatiminput{../final_version/Q4/part2/code/run_q42.py}

\subsubsection*{verify\_q42.py}
\verbatiminput{../final_version/Q4/part2/code/verify_q42.py}

\subsubsection*{check\_information.py}
\verbatiminput{../final_version/Q4/part2/code/check_information.py}

\subsubsection*{report\_q42.py}
\verbatiminput{../final_version/Q4/part2/code/report_q42.py}

\subsection{问题四-3：波动电价下的多时点调整}

\subsubsection*{q3\_data.py}
\verbatiminput{../final_version/Q4/part3/code/q3_data.py}

\subsubsection*{q3\_forecast.py}
\verbatiminput{../final_version/Q4/part3/code/q3_forecast.py}

\subsubsection*{q3\_model.py}
\verbatiminput{../final_version/Q4/part3/code/q3_model.py}

\subsubsection*{run\_q3.py}
\verbatiminput{../final_version/Q4/part3/code/run_q3.py}

\subsubsection*{verify\_q3.py}
\verbatiminput{../final_version/Q4/part3/code/verify_q3.py}

\subsection{公共模块：第二问基础数据与 LP（q2\_base）}

\subsubsection*{q2\_data.py}
\verbatiminput{../final_version/common/q2_base/q2_data.py}

\subsubsection*{run\_1439.py}
\verbatiminput{../final_version/common/q2_base/run_1439.py}

\subsection{公共模块：附件5结果模板表头（主程序依赖）}

\subsubsection*{template\_layout.py}
\verbatiminput{../final_version/common/efficiency/template_layout.py}

\endgroup
